\chapter{Typed computation schemas}
\label{chap:computation-schema}

This chapter defines the types used to state machine and appendix results.  It contains no Lin output, table row, or other computed value.

\section{Finite presentations of an existing Adams page}

% SOURCE: project interface design; this does not assert any real computed value.
\begin{definition}[Finite basis multiplication table]
  \label{def:adams-e2-basis-table}
  \lean{KIP126.AdamsE2.Table}
  \uses{def:classical-adams-e2-e3-shape-slice}
  A table specifies a finite set $D$ of bidegrees containing $(0,0)$,
  dimensions $n(p)$ for $p\in D$, coordinates for the unit, and coefficients
  $c_{p,q,i,j,k}\in\mathbb F_2$ for every product whose source degrees
  $p,q$ and target degree $p+q$ belong to $D$.
  A degree outside $D$ is unknown, not a zero-dimensional space.
\end{definition}

\begin{definition}[Algebra presented by a multiplication table]
  \label{def:adams-e2-table-model}
  \lean{KIP126.AdamsE2.Table.Model,KIP126.AdamsE2.Table.piece}
  \uses{def:adams-e2-basis-table}
  Introduce a generator $X_{p,i}$ of degree $p$ for each recorded basis
  index.  Let $A_T$ be the polynomial algebra over $\mathbb F_2$ modulo
  the unit relation and all covered relations
  \[
    X_{p,i}X_{q,j}=\sum_k c_{p,q,i,j,k}X_{p+q,k}.
  \]
  No relation sets an uncovered product to zero.  The degree-$p$ subspace
  is the image of the weighted homogeneous polynomials of degree $p$.
  This is a model algebra, not a definition of the actual Adams page.
\end{definition}

\begin{definition}[Presentation of the actual $E_2$ page]
  \label{def:adams-e2-table-presentation}
  \lean{KIP126.AdamsE2.PageAlgebra,KIP126.AdamsE2.Presentation,KIP126.AdamsE2.Input}
  \uses{def:adams-e2-table-model,def:classical-adams-e2-e3-shape-slice,def:external-evidence}
  Fix an existing mod--$2$ Adams-shaped spectral sequence $E$ and a
  degree-preserving commutative algebra structure on the direct sum of its
  actual $E_2$ components.  A presentation consists of an algebra map
  $A_T\to E_2$, compatible maps on every homogeneous component, and
  bijectivity on each degree in $D$.  It also supplies bases of the actual
  page components, of the recorded sizes, equal to the images of the
  corresponding model generators.  Existence of this presentation is
  supplied as explicit external evidence with provenance; reading the
  numerical table does not prove it.
\end{definition}

\begin{lemma}[Transporting a table product to the page]
  \label{lem:adams-e2-table-product}
  \lean{KIP126.AdamsE2.Table.generator_mul,KIP126.AdamsE2.Presentation.basis_mul}
  \uses{def:adams-e2-table-presentation}
  Under a supplied presentation, every covered table relation is an
  equality for the actual $E_2$ multiplication on the specified basis.
\end{lemma}
\begin{proof}
  The polynomial relation is zero in the quotient model.  Apply the
  presentation's algebra homomorphism and use injectivity of the inclusion
  of the target component into the direct sum.
\end{proof}

\begin{definition}[The classes obtained from an imported table]
  \label{def:adams-e2-imported-h6}
  \lean{KIP126.AdamsE2.Input.h6,KIP126.AdamsE2.Input.h6Square}
  \uses{def:adams-e2-table-presentation}
  If the table contains $(1,64)$ with dimension one, define $h_6$ as the
  image of its sole basis vector.  Define $h_6^2$ using the actual page
  multiplication.  Its nonvanishing or permanence is not a field of this
  definition.  This interface alone does not identify an arbitrary
  Adams-shaped sequence with the sphere's sequence.
\end{definition}

\begin{lemma}[Three-cell illustration]
  \label{lem:adams-e2-three-cell-example}
  \lean{KIP126.Examples.AdamsE2Table.x_mul_x,KIP126.Examples.AdamsE2Table.y_ne_zero,KIP126.Examples.AdamsE2Table.x_mul_y_ne_zero,KIP126.Examples.AdamsE2Table.imported_h6_square_ne_zero}
  \uses{def:adams-e2-imported-h6,lem:adams-e2-table-product}
  Consider the illustrative table with one-dimensional cells at $(0,0)$,
  $(1,64)$ and $(2,128)$, with unit and $x^2=y$.  In its model, $y\ne0$
  and the uncovered product $xy\ne0$.  For any supplied presentation of
  this table on an existing page, the resulting $h_6^2$ is nonzero.
  This is not a claim that the illustration is actual sphere Adams data.
\end{lemma}
\begin{proof}
  Map the model to $\mathbb F_2[X]$ by $x\mapsto X$, $y\mapsto X^2$
  and the unit generator to $1$.  All table relations vanish, whereas
  $X^2$ and $X^3$ are nonzero.  On the actual page, the table product is
  the specified basis vector in degree $(2,128)$ and is therefore nonzero.
\end{proof}

\section{Typed catalogue schemas}

% SOURCE: aimpaper/main.tex:1930-1977, Example exam:Mahowald.
\begin{definition}[Atiyah--Hirzebruch cell notation for $S^0/\nu_{\mathrm{Hopf}}$]
  \label{def:hopf-cofiber-cell-notation}
  \uses{def:cofiber-triangle,def:steenrod-ext}
  \notready
  In the Ext long exact sequence of
  $S^3\xrightarrow{\nu_{\mathrm{Hopf}}}S^0\to S^0/\nu_{\mathrm{Hopf}}\to S^4$, write
  $a[4]$ for a class restricting nontrivially to the top cell as $a$, and
  $b[0]$ for a class induced from the bottom-cell class $b$.  The notation is
  typed by spectrum, cell, bidegree, and a chosen dataset identifier.
\end{definition}

% SOURCE: aimpaper/main.tex:2738-3290, all twelve tables.
\begin{definition}[Appendix table identifier]
  \label{def:appendix-table-id}
  \leanok
  \lean{KIP126.Computation.AppendixTableId}
  \lean{KIP126.Computation.AppendixTableId.paperNumber}
  \lean{KIP126.Computation.AppendixTableId.texLabel}
  \lean{KIP126.Computation.AppendixTableId.sourceStart}
  \lean{KIP126.Computation.AppendixTableId.sourceEnd}
  \lean{KIP126.Computation.AppendixTableId.pdfPage}
  \lean{KIP126.Computation.AppendixTableId.spectrum}
  \lean{KIP126.Computation.AppendixTableId.stem}
  \lean{KIP126.Computation.AppendixTableId.filtrationRange}
  The table identifier is a twelve-element finite type containing
  $C\nu126$, $S122$, $S123$, the two $S124$ bands, the two $S125$ bands, the
  two $S126$ bands, and the three $S127$ bands.  It records the paper table
  number, spectrum, stem, filtration interval, source line range, and PDF page.
\end{definition}

% SOURCE: aimpaper/main.tex:2738-3290, class-expression column; aimpaper/112.tex:8-1026.
\begin{definition}[Typed Ext class expression]
  \label{def:typed-class-expression}
  \uses{def:steenrod-ext,def:hopf-cofiber-cell-notation}
  \leanok
  \lean{KIP126.Computation.DatasetVersion}
  \lean{KIP126.Computation.Cell}
  \lean{KIP126.Computation.ClassAtom}
  \lean{KIP126.Computation.ClassTerm}
  \lean{KIP126.Computation.ExpressionMetadata}
  \lean{KIP126.Computation.TypedClassExpression}
  \lean{KIP126.Computation.ClassTerm.validBool}
  \lean{KIP126.Computation.ClassTerm.Valid}
  \lean{KIP126.Computation.TypedClassExpression.DegreeValid}
  \lean{KIP126.Computation.TypedClassExpression.Valid}
  A typed class expression is a sparse $\mathbb F_2$ expression built from a
  dataset-versioned generator $x_{n,s,k}$, a named class, product, and finite
  sum, together with spectrum, stem, filtration, internal degree, and optional
  bottom- or top-cell tag.  The Lean data records degree validity as a
  predicate; a finite candidate set represents ambiguity without converting it
  to free-form text.
\end{definition}

% SOURCE: aimpaper/main.tex:2738-3290, differential/status columns.
\begin{definition}[Appendix row status]
  \label{def:appendix-row-status}
  \uses{def:typed-class-expression}
  \leanok
  \lean{KIP126.Computation.DifferentialDirection}
  \lean{KIP126.Computation.AppendixRowStatus}
  \lean{KIP126.Computation.AppendixRowStatus.isIncoming}
  \lean{KIP126.Computation.AppendixRowStatus.unresolved}
  Row status is one of:
  \begin{enumerate}
    \item an explicit zero filtration band;
    \item a permanent class;
    \item a known outgoing $d_r$ with a nonempty finite target set;
    \item a known incoming $d_r$ with a nonempty finite source set;
    \item survival through a stated page with the next outgoing differential
      unresolved.
  \end{enumerate}
  An incoming marker $d_r^{-1}$ is not an inverse map.  A printed $d_r?$ does
  not assert existence of a nonzero differential.  A ``possibly'' term is a
  finite target ambiguity attached to a known relation.
\end{definition}

% SOURCE: aimpaper/main.tex:2738-3290; duplicated source/target presentations across adjacent stems.
\begin{definition}[Differential relation identifier]
  \label{def:differential-relation-id}
  \uses{def:appendix-row-status}
  \leanok
  \lean{KIP126.Computation.DifferentialRelationId}
  \lean{KIP126.Computation.DifferentialRelationId.Valid}
  \lean{KIP126.Computation.DifferentialRelationId.validBool}
  Each known differential has a stable relation identifier, length, typed
  source expression, and typed target candidate set.  The row schema carries
  this identifier explicitly; joining outgoing and incoming presentations is
  a later validation step.
\end{definition}

% SOURCE: aimpaper/main.tex:2738-3290.
\begin{definition}[Appendix row catalogue]
  \label{def:appendix-row-catalogue}
  \uses{def:appendix-table-id,def:typed-class-expression,def:appendix-row-status,def:differential-relation-id}
  \leanok
  \lean{KIP126.Computation.AppendixRowLocator}
  \lean{KIP126.Computation.AppendixRow}
  \lean{KIP126.Computation.AppendixRow.KeyValid}
  \lean{KIP126.Computation.AppendixRow.MetadataValid}
  \lean{KIP126.Computation.AppendixRow.RelationValid}
  \lean{KIP126.Computation.AppendixRow.Valid}
  \lean{KIP126.Computation.AppendixRow.validBool}
  \lean{KIP126.Computation.AppendixZeroBand}
  \lean{KIP126.Computation.AppendixZeroBand.Valid}
  \lean{KIP126.Computation.AppendixZeroBand.validBool}
  \lean{KIP126.Computation.expressionAt}
  \lean{KIP126.Computation.sourceExpression}
  \lean{KIP126.Computation.appendixRows}
  \lean{KIP126.Computation.appendixRows_length}
  \lean{KIP126.Computation.appendixRows_keys_nodup}
  \lean{KIP126.Computation.appendixRows_valid}
  \lean{KIP126.Computation.appendixZeroBands}
  \lean{KIP126.Computation.appendixZeroBands_length}
  \lean{KIP126.Computation.appendixZeroBands_valid}
  The source-shaped catalogue contains all 401 nonempty rows printed in the
  twelve tables, with table key, stem, filtration, typed expression, status,
  relation key, and exact source line.  The nine explicitly empty filtration
  bands are represented separately as display records.
\end{definition}

% SOURCE: LWXMachine, Section 2.1 and Tables 1--5; Zenodo record 14875701,
% version v126.3.cw49; PROJECT_BOUNDARY.md, Appendix data.
\begin{definition}[Lin CW-spectrum catalogue]
  \label{def:lin-spectrum-catalogue}
  \uses{def:source-ref}
  \notready
  The computation input contains a finite, versioned catalogue of the 49 CW
  spectra used by Lin--Wang--Xu.  Each entry records its stable name, cell
  presentation and suspensions, whether it is used as a ring or module
  spectrum, its maximum computed internal degree, and source locators for the
  corresponding files in the Zenodo archive.  Spectrum names appearing in a
  table, map, or proof record are typed references to this catalogue rather
  than free-form strings.
\end{definition}

% SOURCE: LWXMachine, Section 2.1, Notations 2.1--2.5 and the
% *_AdamsE2_generators/relations/basis files described there.
\begin{definition}[Lin Adams $E_2$-page catalogue]
  \label{def:lin-e2-page-catalogue}
  \uses{def:lin-spectrum-catalogue,def:typed-class-expression,def:external-evidence}
  \notready
  For every spectrum in the Lin catalogue, an $E_2$-page record contains the
  complete finite generator, relation, and basis tables through its declared
  internal-degree bound.  It includes typed multiplication or module-action
  tables only when the spectrum catalogue supplies the corresponding ring or
  module structure.  An external-evidence value identifies the three archived
  files and proves that their normalized serialization is the value stored in
  the record.
\end{definition}

% SOURCE: LWXMachine, Section 2.2 and Tables 6--10; the map_*.csv files in
% Zenodo record 14875701.
\begin{definition}[Lin map catalogue]
  \label{def:lin-map-catalogue}
  \uses{def:lin-spectrum-catalogue,def:lin-e2-page-catalogue,def:external-evidence}
  \notready
  The map catalogue contains the 180 versioned maps used by the computation.
  A record gives source and target spectrum identifiers, suspension degree,
  geometric map kind, maximum computed degree, and the induced linear map on
  the stored $E_2$ bases.  Cofiber-sequence membership and composition keys are
  explicit fields, and the evidence locator names the archived map table rather
  than merely the display name of the map.
\end{definition}

% SOURCE: LWXMachine, Section 2.3 and Table 11; the d_2 columns in the
% per-spectrum Adams E_2 basis files;
% aimpaper/main.tex:2785--2792 for the additional manual inputs.
\begin{definition}[Lin initial $d_2$ catalogue]
  \label{def:lin-d2-catalogue}
  \uses{def:lin-e2-page-catalogue,def:differential-relation-id,def:external-evidence}
  \notready
  For every spectrum and degree range declared by the computation, the initial
  $d_2$ catalogue records all supplied $d_2$ values as typed differential
  relations and carries a range-completeness certificate.  A missing row is
  therefore distinguishable from a certified zero differential.  Later manual
  $d_3$, $d_5$, and $d_6$ inputs are stored separately and are never
  reconstructed from program output.
\end{definition}

% SOURCE: LWXMachine, Section 3, especially the Adams-SS, cofiber-SS, and
% Tables-of-Proofs formats; proofs_csv.rar in Zenodo record 14875701.
\begin{definition}[Lin propagated-output record]
  \label{def:lin-propagated-output-record}
  \uses{def:lin-spectrum-catalogue,def:lin-map-catalogue,def:lin-d2-catalogue,def:differential-relation-id,def:external-evidence}
  \notready
  A propagated-output record is one of a proved Adams differential, a proved
  map extension, a disproof of a candidate value, or an explicitly unresolved
  finite candidate set.  It stores the source inputs, propagation rule,
  normalized source and target, page or extension length, proof-database key,
  program version, and archived output row.  A disproof record contains the
  rejected candidate and contradiction key; it is not represented as an absent
  differential row.
\end{definition}

% SOURCE: LWXMachine Sections 2--3; Zenodo record 14875701, version
% v126.3.cw49; PROJECT_BOUNDARY.md, Appendix data.
\begin{definition}[Lin computation bundle]
  \label{def:lin-computation-bundle}
  \uses{def:lin-spectrum-catalogue,def:lin-e2-page-catalogue,def:lin-map-catalogue,def:lin-d2-catalogue,def:lin-propagated-output-record}
  \notready
  A Lin computation bundle contains the five catalogues above with one common
  archive version and hashes for the program, database, CSV export, and proof
  files.  Its validation data states that every map references stored source
  and target pages, every propagated output references stored inputs, and no
  two records assign incompatible values to the same typed relation key.
\end{definition}

% SOURCE: PROJECT_BOUNDARY.md, Appendix data; aimpaper/main.tex:2785-2800.
\begin{definition}[Appendix evidence record]
  \label{def:appendix-evidence-record}
  \uses{def:appendix-table-id,def:appendix-row-status,def:differential-relation-id,def:external-evidence}
  \notready
  An appendix evidence record contains a table identifier, a unique row key,
  spectrum and bidegree, typed class expression or zero band, row status,
  relation identifier when applicable, paper locator, computation-artifact
  locator, program version, input hash, method, and an explicit evidence proof
  of the proposition represented by that status.
\end{definition}

\section{Input bundles and output contracts}

% SOURCE: PROJECT_BOUNDARY.md, Appendix data and External inputs;
% aimpaper/main.tex:2738--3290.
\begin{definition}[Near-$126$ computation input]
  \label{def:near126-computation-input}
  \uses{def:lin-computation-bundle,def:appendix-evidence-record,def:external-result,def:external-evidence}
  \notready
  A near-$126$ computation input is an explicit record containing a Lin
  computation bundle, the twelve finite appendix tables, the nine zero-band
  records, the three manual differential inputs, the $tmf$-detection evidence,
  and the finite high-stem product and indeterminacy evidence consumed in
  Section~7.  The latter are named fields, including
  \texttt{theta5TorsionRanges}, \texttt{theta5SquareTmf},
  \texttt{todaCandidateProducts},
  \texttt{twoExtensionIndeterminacyProducts},
  \texttt{hopfLiftObstructions}, \texttt{stem122Products}, and
  \texttt{cnuIncomingRange}; a proof cannot request an anonymous
  \texttt{ExternalEvidence}.  Every field is an external-evidence or
  external-result value, as
  appropriate, with a source locator; project theorems receive this record as a
  parameter and prove all finite coverage and negative-search conclusions from
  its fields.
\end{definition}
